\begin{abcliste}
\abc Compare the kinetic energy with the rest energy $E_0 = m c^2$. The classical formula is good only if $E_\mathrm{kin}$ is small compared with $E_0$ (a few per cent at most).

Electron:
\begin{align}
 E_0 &= \EzeF = \Eze \approx \EzeP{3}{3} \\
 \frac{E_\mathrm{kin}}{E_0} &= \frac{\EeO}{\EzeP{3}{3}} \approx \xeP{0}{2}
\end{align}
That is about \SI{20}{\percent}: for the electron, \result{\text{relativity is needed}}.

Proton:
\begin{align}
 E_0 &= \EzpF = \Ezp \approx \EzpP{6}{3} \\
 \frac{E_\mathrm{kin}}{E_0} &= \frac{\EpO}{\EzpP{6}{3}} \approx \xpP{0}{2}
\end{align}
Only about \SI{2}{\percent}: for the proton, \result{\text{the classical formula is good enough}}.

\abc The rest energy of one atomic mass unit is
\begin{align}
 m_u c^2 &= \EuF = \Eu \approx \EuP{9}{4}
\end{align}
so \SI{1}{\atomicmassunit} $= \EuP{9}{4}/c^2$ and
\begin{align}
 m &= \ECF = 12\times\EuP{9}{4}/c^2 \approx \result{\ECP{9}{3}/c^2}
\end{align}
(Dividing by \EuP{9}{4} instead of multiplying, or forgetting the prefix, are the typical slips.)
\end{abcliste}
